\documentclass[11pt,a4paper]{article}

\usepackage[utf8]{inputenc}
\usepackage[T1]{fontenc}
\usepackage{amsmath,amssymb,amsthm}
\usepackage{algorithm}
\usepackage{algpseudocode}
\usepackage{graphicx}
\usepackage{hyperref}
\usepackage{listings}
\usepackage{xcolor}
\usepackage{booktabs}
\usepackage{geometry}
\usepackage{tikz}
\usetikzlibrary{arrows.meta,positioning,shapes}

\geometry{margin=1in}

\theoremstyle{definition}
\newtheorem{definition}{Definition}
\newtheorem{theorem}{Theorem}
\newtheorem{lemma}{Lemma}
\newtheorem{property}{Property}

\lstset{
    basicstyle=\ttfamily\small,
    keywordstyle=\color{blue},
    commentstyle=\color{gray},
    stringstyle=\color{red},
    breaklines=true,
    frame=single
}

\title{VRF Precompile for Verifiable Randomness:\\Chainlink-Compatible On-Chain Random Number Generation}

\author{Zach Kelling\\
\textit{Hanzo Industries}\\
\textit{Lux Industries}\\
\textit{Zoo Labs Foundation}\\
\texttt{zach@hanzo.ai}}

\date{2025}

\begin{document}

\maketitle

\begin{abstract}
We present a precompiled contract specification for Verifiable Random Functions (VRFs), enabling provably fair on-chain randomness generation. VRFs provide cryptographic proofs that random outputs were correctly computed from inputs, preventing manipulation by oracle operators or miners. This paper specifies a precompile at address \texttt{0x19} implementing ECVRF over secp256k1 (Chainlink-compatible) and ed25519 (IETF standard), with extensions for BLS-based threshold VRFs. We analyze security properties, gas models, and integration patterns for gaming, NFT minting, validator selection, and decentralized lotteries.
\end{abstract}

\section{Introduction}

Secure randomness is fundamental to blockchain applications yet notoriously difficult to achieve. Block hashes are manipulable by miners, commit-reveal schemes require multiple transactions, and external oracles introduce trust assumptions. Verifiable Random Functions solve this by producing pseudorandom outputs with publicly verifiable proofs.

\subsection{The Randomness Problem}

\begin{enumerate}
    \item \textbf{Unpredictability}: Outputs must be computationally indistinguishable from random
    \item \textbf{Unbiasability}: No party can influence the output
    \item \textbf{Verifiability}: Anyone can verify correct computation
    \item \textbf{Uniqueness}: Each input maps to exactly one valid output
\end{enumerate}

\subsection{VRF Solution}

A VRF is a keyed function where:
\begin{itemize}
    \item Only the secret key holder can compute the output
    \item The output is deterministic (same input $\Rightarrow$ same output)
    \item A proof convinces anyone of correct computation
    \item Without the key, outputs appear random
\end{itemize}

\subsection{Contributions}

\begin{itemize}
    \item VRF precompile specification at address \texttt{0x19}
    \item Support for ECVRF-secp256k1 (Chainlink-compatible)
    \item Support for ECVRF-ed25519 (IETF RFC 9381)
    \item BLS-based distributed VRF for threshold randomness
    \item Gas-efficient verification suitable for on-chain use
\end{itemize}

\section{Preliminaries}

\subsection{Verifiable Random Functions}

\begin{definition}[VRF]
A Verifiable Random Function consists of three algorithms:
\begin{itemize}
    \item $\text{KeyGen}() \to (sk, pk)$: Generate key pair
    \item $\text{Prove}(sk, \alpha) \to (\beta, \pi)$: Compute output $\beta$ and proof $\pi$ for input $\alpha$
    \item $\text{Verify}(pk, \alpha, \beta, \pi) \to \{0, 1\}$: Verify proof
\end{itemize}
\end{definition}

\subsection{Security Properties}

\begin{property}[Uniqueness]
For any $pk$, $\alpha$, there exists at most one $\beta$ for which a valid proof exists.
\end{property}

\begin{property}[Pseudorandomness]
Without $sk$, the output $\beta$ is computationally indistinguishable from random, even given the public key and other input-output pairs.
\end{property}

\begin{property}[Provability]
If $\beta = \text{VRF}_{sk}(\alpha)$, then $\text{Verify}(pk, \alpha, \beta, \pi) = 1$.
\end{property}

\subsection{ECVRF Construction}

ECVRF uses elliptic curve cryptography:

\begin{enumerate}
    \item Hash input to curve point: $H = \text{HashToCurve}(\alpha)$
    \item Compute VRF output point: $\Gamma = sk \cdot H$
    \item Generate Schnorr-like proof: $\pi = (c, s)$
    \item Hash $\Gamma$ to produce final output: $\beta = \text{Hash}(\Gamma)$
\end{enumerate}

\section{ECVRF Specification}

\subsection{ECVRF-secp256k1 (Chainlink Compatible)}

This variant ensures compatibility with Chainlink VRF v2.

\subsubsection{Parameters}

\begin{itemize}
    \item Curve: secp256k1
    \item Hash: keccak256
    \item Suite: \texttt{secp256k1\_keccak256\_tai}
\end{itemize}

\subsubsection{Hash to Curve}

\begin{algorithm}
\caption{HashToCurve (Try-and-Increment)}
\begin{algorithmic}[1]
\Require Input $\alpha$, public key $pk$
\Ensure Point $H \in E$
\State $\text{ctr} \gets 0$
\While{$\text{ctr} < 256$}
    \State $h \gets \text{keccak256}(\text{suite} \| pk \| \alpha \| \text{ctr})$
    \State $x \gets h[0:32] \mod p$
    \State $y^2 \gets x^3 + 7 \mod p$
    \If{$y^2$ is quadratic residue}
        \State $y \gets \sqrt{y^2}$
        \If{$y$ is odd and $h[31] \& 1 = 0$}
            \State $y \gets p - y$
        \EndIf
        \Return $(x, y)$
    \EndIf
    \State $\text{ctr} \gets \text{ctr} + 1$
\EndWhile
\State \textbf{fail}
\end{algorithmic}
\end{algorithm}

\subsubsection{Prove Algorithm}

\begin{algorithm}
\caption{ECVRF Prove}
\begin{algorithmic}[1]
\Require Secret key $sk$, input $\alpha$
\Ensure Output $\beta$, proof $\pi$
\State $pk \gets sk \cdot G$
\State $H \gets \text{HashToCurve}(\alpha, pk)$
\State $\Gamma \gets sk \cdot H$
\State $k \gets \text{NonceGeneration}(sk, H)$
\State $U \gets k \cdot G$
\State $V \gets k \cdot H$
\State $c \gets \text{HashPoints}(pk, H, \Gamma, U, V) \mod q$
\State $s \gets (k + c \cdot sk) \mod q$
\State $\beta \gets \text{keccak256}(\Gamma_x \| \Gamma_y)$
\Return $(\beta, (\Gamma, c, s))$
\end{algorithmic}
\end{algorithm}

\subsubsection{Verify Algorithm}

\begin{algorithm}
\caption{ECVRF Verify}
\begin{algorithmic}[1]
\Require Public key $pk$, input $\alpha$, output $\beta$, proof $\pi = (\Gamma, c, s)$
\Ensure Valid or Invalid
\State $H \gets \text{HashToCurve}(\alpha, pk)$
\State $U \gets s \cdot G - c \cdot pk$
\State $V \gets s \cdot H - c \cdot \Gamma$
\State $c' \gets \text{HashPoints}(pk, H, \Gamma, U, V) \mod q$
\If{$c \neq c'$}
    \Return Invalid
\EndIf
\State $\beta' \gets \text{keccak256}(\Gamma_x \| \Gamma_y)$
\If{$\beta \neq \beta'$}
    \Return Invalid
\EndIf
\Return Valid
\end{algorithmic}
\end{algorithm}

\subsection{ECVRF-ed25519 (IETF RFC 9381)}

\subsubsection{Parameters}

\begin{itemize}
    \item Curve: ed25519 (Curve25519)
    \item Hash: SHA-512
    \item Suite: \texttt{ECVRF-ED25519-SHA512-ELL2}
\end{itemize}

\subsubsection{Encoding}

\begin{itemize}
    \item Public key: 32 bytes (compressed)
    \item Proof: 80 bytes ($\Gamma$ + $c$ + $s$)
    \item Output: 64 bytes (before truncation)
\end{itemize}

\section{Precompile Specification}

\subsection{Contract Address}

The VRF precompile occupies address \texttt{0x19}.

\subsection{Input Format}

\begin{lstlisting}
struct VRFInput {
    uint8 suite;        // VRF suite identifier
    bytes publicKey;    // Signer's public key
    bytes alpha;        // VRF input (seed)
    bytes pi;           // VRF proof
}
\end{lstlisting}

Suite identifiers:
\begin{table}[h]
\centering
\begin{tabular}{cl}
\toprule
\textbf{ID} & \textbf{Suite} \\
\midrule
0x01 & ECVRF-secp256k1-keccak256 (Chainlink) \\
0x02 & ECVRF-ed25519-SHA512 (IETF) \\
0x03 & BLS-VRF-BLS12381 (Threshold) \\
\bottomrule
\end{tabular}
\caption{VRF suite identifiers}
\label{tab:suites}
\end{table}

\subsection{Output Format}

Returns 64 bytes:
\begin{itemize}
    \item Bytes 0-31: Verification result (1 = valid, 0 = invalid)
    \item Bytes 32-63: VRF output $\beta$ (if valid)
\end{itemize}

\subsection{Gas Cost Model}

\begin{equation}
\text{gas} = G_{\text{base}} + G_{\text{suite}} + G_{\text{alpha}} \cdot \lceil |\alpha| / 32 \rceil
\end{equation}

\begin{table}[h]
\centering
\begin{tabular}{lcc}
\toprule
\textbf{Suite} & $G_{\text{suite}}$ & \textbf{Total (32-byte input)} \\
\midrule
ECVRF-secp256k1 & 25,000 & 26,500 \\
ECVRF-ed25519 & 18,000 & 19,500 \\
BLS-VRF & 180,000 & 181,500 \\
\bottomrule
\end{tabular}
\caption{VRF gas costs ($G_{\text{base}} = 1500$, $G_{\text{alpha}} = 12$)}
\label{tab:gas}
\end{table}

\section{BLS-Based Threshold VRF}

\subsection{Motivation}

Single-key VRFs centralize trust in the key holder. Threshold VRFs distribute this trust across $n$ parties, requiring $t$ to cooperate for output generation.

\subsection{DVRF Construction}

\begin{definition}[Distributed VRF]
A $(t, n)$-threshold VRF satisfies:
\begin{enumerate}
    \item Any $t$ parties can compute the VRF output
    \item Fewer than $t$ parties learn nothing about the output
    \item Output is unique regardless of which $t$ parties participate
\end{enumerate}
\end{definition}

\subsection{BLS-DVRF Protocol}

\begin{algorithm}
\caption{BLS-DVRF Partial Evaluation}
\begin{algorithmic}[1]
\Require Party $i$'s secret share $sk_i$, input $\alpha$
\Ensure Partial output $\sigma_i$
\State $H \gets \text{HashToG1}(\alpha)$
\State $\sigma_i \gets sk_i \cdot H$
\Return $\sigma_i$
\end{algorithmic}
\end{algorithm}

\begin{algorithm}
\caption{BLS-DVRF Combine}
\begin{algorithmic}[1]
\Require Partial outputs $\{(\sigma_i, i)\}$ for $|S| \geq t$ parties
\Ensure Combined output $\beta$, proof $\pi$
\State Compute Lagrange coefficients $\{\lambda_i\}_{i \in S}$
\State $\Gamma \gets \sum_{i \in S} \lambda_i \cdot \sigma_i$
\State $\beta \gets \text{Hash}(\Gamma)$
\State $\pi \gets \Gamma$
\Return $(\beta, \pi)$
\end{algorithmic}
\end{algorithm}

\begin{algorithm}
\caption{BLS-DVRF Verify}
\begin{algorithmic}[1]
\Require Aggregate public key $pk$, input $\alpha$, output $\beta$, proof $\Gamma$
\Ensure Valid or Invalid
\State $H \gets \text{HashToG1}(\alpha)$
\If{$e(\Gamma, G_2) \neq e(H, pk)$}
    \Return Invalid
\EndIf
\State $\beta' \gets \text{Hash}(\Gamma)$
\If{$\beta \neq \beta'$}
    \Return Invalid
\EndIf
\Return Valid
\end{algorithmic}
\end{algorithm}

\section{Security Analysis}

\subsection{Uniqueness}

\begin{theorem}[VRF Uniqueness]
For ECVRF, the output $\beta$ is uniquely determined by $(pk, \alpha)$ under the Computational Diffie-Hellman assumption.
\end{theorem}

\begin{proof}
Suppose two valid outputs $\beta_1 \neq \beta_2$ exist for the same $(pk, \alpha)$. Then there exist $\Gamma_1 \neq \Gamma_2$ with valid proofs. But both satisfy $\Gamma = sk \cdot H$ for the same $H$, implying $\Gamma_1 = \Gamma_2$, contradiction.
\end{proof}

\subsection{Pseudorandomness}

\begin{theorem}[VRF Pseudorandomness]
ECVRF outputs are pseudorandom under the Decisional Diffie-Hellman assumption.
\end{theorem}

The hash of $\Gamma = sk \cdot H$ is indistinguishable from random because $\Gamma$ is a DDH tuple.

\subsection{Unpredictability Window}

Even if an adversary knows future inputs (e.g., block numbers), they cannot predict outputs without the secret key.

\subsection{Front-Running Resistance}

VRF outputs cannot be predicted before the proof is published, preventing front-running in applications like NFT reveals.

\section{Integration Patterns}

\subsection{Chainlink VRF Compatible}

\begin{lstlisting}[language=Solidity]
contract VRFConsumer {
    address constant VRF_PRECOMPILE = address(0x19);

    bytes public vrfPublicKey;
    mapping(uint256 => bytes32) public randomResults;

    event RandomnessRequested(uint256 indexed requestId, bytes32 seed);
    event RandomnessFulfilled(uint256 indexed requestId, bytes32 randomness);

    function requestRandomness(
        uint256 requestId
    ) external returns (bytes32 seed) {
        seed = keccak256(abi.encodePacked(
            block.chainid,
            address(this),
            requestId,
            block.number
        ));

        emit RandomnessRequested(requestId, seed);
        return seed;
    }

    function fulfillRandomness(
        uint256 requestId,
        bytes32 seed,
        bytes calldata proof
    ) external {
        // Verify VRF proof
        (bool success, bytes memory result) = VRF_PRECOMPILE.staticcall(
            abi.encodePacked(
                uint8(0x01), // ECVRF-secp256k1
                vrfPublicKey,
                seed,
                proof
            )
        );
        require(success, "VRF call failed");

        // Parse result
        uint256 valid = uint256(bytes32(result[0:32]));
        require(valid == 1, "Invalid VRF proof");

        bytes32 randomness = bytes32(result[32:64]);
        randomResults[requestId] = randomness;

        emit RandomnessFulfilled(requestId, randomness);
    }

    function expandRandomness(
        bytes32 randomness,
        uint256 count
    ) public pure returns (uint256[] memory) {
        uint256[] memory expanded = new uint256[](count);
        for (uint256 i = 0; i < count; i++) {
            expanded[i] = uint256(keccak256(abi.encodePacked(randomness, i)));
        }
        return expanded;
    }
}
\end{lstlisting}

\subsection{NFT Reveal}

\begin{lstlisting}[language=Solidity]
contract RevealableNFT {
    address constant VRF = address(0x19);

    struct Collection {
        bytes vrfPublicKey;
        uint256 totalSupply;
        uint256 revealBlock;
        bytes32 revealSeed;
        bytes32 randomness;
        bool revealed;
    }

    mapping(uint256 => Collection) public collections;
    mapping(uint256 => mapping(uint256 => uint256)) public tokenToMetadata;

    function commitReveal(
        uint256 collectionId
    ) external {
        Collection storage c = collections[collectionId];
        require(!c.revealed, "Already revealed");

        c.revealBlock = block.number + 10; // 10 block delay
        c.revealSeed = keccak256(abi.encodePacked(
            collectionId,
            block.number,
            blockhash(block.number - 1)
        ));
    }

    function executeReveal(
        uint256 collectionId,
        bytes calldata vrfProof
    ) external {
        Collection storage c = collections[collectionId];
        require(block.number >= c.revealBlock, "Too early");
        require(!c.revealed, "Already revealed");

        // Include blockhash in seed for additional entropy
        bytes32 finalSeed = keccak256(abi.encodePacked(
            c.revealSeed,
            blockhash(c.revealBlock)
        ));

        // Verify VRF
        (bool success, bytes memory result) = VRF.staticcall(
            abi.encodePacked(
                uint8(0x01),
                c.vrfPublicKey,
                finalSeed,
                vrfProof
            )
        );
        require(success && uint256(bytes32(result[0:32])) == 1, "Invalid VRF");

        c.randomness = bytes32(result[32:64]);
        c.revealed = true;

        // Shuffle metadata assignment
        assignMetadata(collectionId);
    }

    function assignMetadata(uint256 collectionId) internal {
        Collection storage c = collections[collectionId];

        // Fisher-Yates shuffle using VRF randomness
        for (uint256 i = 0; i < c.totalSupply; i++) {
            uint256 j = i + uint256(keccak256(abi.encodePacked(
                c.randomness, i
            ))) % (c.totalSupply - i);

            // Swap
            uint256 temp = tokenToMetadata[collectionId][i];
            if (temp == 0) temp = i;

            uint256 other = tokenToMetadata[collectionId][j];
            if (other == 0) other = j;

            tokenToMetadata[collectionId][i] = other;
            tokenToMetadata[collectionId][j] = temp;
        }
    }
}
\end{lstlisting}

\subsection{Validator Selection}

\begin{lstlisting}[language=Solidity]
contract ValidatorSelection {
    address constant VRF = address(0x19);

    struct Validator {
        address addr;
        uint256 stake;
        bytes vrfPublicKey;
    }

    Validator[] public validators;
    uint256 public totalStake;

    function selectProposer(
        uint256 epoch,
        uint256 slot,
        bytes calldata vrfProof
    ) external view returns (address) {
        // Seed from epoch and slot
        bytes32 seed = keccak256(abi.encodePacked(epoch, slot));

        // Expected proposer submits their VRF proof
        uint256 proposerIndex = slot % validators.length;
        Validator storage v = validators[proposerIndex];

        // Verify VRF
        (bool success, bytes memory result) = VRF.staticcall(
            abi.encodePacked(
                uint8(0x01),
                v.vrfPublicKey,
                seed,
                vrfProof
            )
        );
        require(success && uint256(bytes32(result[0:32])) == 1, "Invalid VRF");

        // Use VRF output to determine if proposer is selected
        // (stake-weighted sortition)
        bytes32 randomness = bytes32(result[32:64]);
        uint256 threshold = v.stake * type(uint256).max / totalStake;

        if (uint256(randomness) < threshold) {
            return v.addr;
        }

        revert("Proposer not selected");
    }
}
\end{lstlisting}

\subsection{Decentralized Lottery}

\begin{lstlisting}[language=Solidity]
contract DecentralizedLottery {
    address constant VRF = address(0x19);

    struct Lottery {
        uint256 endBlock;
        uint256 ticketPrice;
        uint256 prizePool;
        address[] participants;
        bytes oraclePublicKey;
        address winner;
        bool finalized;
    }

    mapping(uint256 => Lottery) public lotteries;

    function buyTicket(uint256 lotteryId) external payable {
        Lottery storage l = lotteries[lotteryId];
        require(block.number < l.endBlock, "Lottery ended");
        require(msg.value == l.ticketPrice, "Wrong price");

        l.participants.push(msg.sender);
        l.prizePool += msg.value;
    }

    function finalizeLottery(
        uint256 lotteryId,
        bytes calldata vrfProof
    ) external {
        Lottery storage l = lotteries[lotteryId];
        require(block.number >= l.endBlock, "Not ended");
        require(!l.finalized, "Already finalized");
        require(l.participants.length > 0, "No participants");

        // Seed includes blockhash for unpredictability
        bytes32 seed = keccak256(abi.encodePacked(
            lotteryId,
            l.endBlock,
            blockhash(l.endBlock)
        ));

        // Verify VRF
        (bool success, bytes memory result) = VRF.staticcall(
            abi.encodePacked(
                uint8(0x01),
                l.oraclePublicKey,
                seed,
                vrfProof
            )
        );
        require(success && uint256(bytes32(result[0:32])) == 1, "Invalid VRF");

        bytes32 randomness = bytes32(result[32:64]);
        uint256 winnerIndex = uint256(randomness) % l.participants.length;

        l.winner = l.participants[winnerIndex];
        l.finalized = true;

        // Transfer prize
        payable(l.winner).transfer(l.prizePool);
    }
}
\end{lstlisting}

\section{Comparison with Alternatives}

\begin{table}[h]
\centering
\begin{tabular}{lcccc}
\toprule
\textbf{Method} & \textbf{Unpredictable} & \textbf{Unbiasable} & \textbf{Verifiable} & \textbf{Gas} \\
\midrule
Block hash & No & No & Yes & 20 \\
Commit-reveal & Yes & Partial & Yes & 40,000+ \\
Chainlink VRF & Yes & Yes & Yes & 200,000+ \\
VRF Precompile & Yes & Yes & Yes & 26,500 \\
\bottomrule
\end{tabular}
\caption{Randomness methods comparison}
\label{tab:compare}
\end{table}

\section{Implementation Considerations}

\subsection{Determinism}

The precompile must produce identical results across all nodes:
\begin{itemize}
    \item Constant-time operations for security
    \item Deterministic hash-to-curve
    \item No floating-point arithmetic
\end{itemize}

\subsection{Input Validation}

\begin{itemize}
    \item Public key must be on curve
    \item Proof components must be in valid range
    \item Input length must match suite requirements
\end{itemize}

\subsection{Error Handling}

Invalid inputs return zeros (not revert) for gas predictability:
\begin{itemize}
    \item Invalid public key: Return (0, 0)
    \item Invalid proof: Return (0, 0)
    \item Verification failure: Return (0, $\beta$) where $\beta$ is computed but invalid
\end{itemize}

\section{Related Work}

\subsection{Chainlink VRF}

Chainlink VRF v2 uses ECVRF-secp256k1 with:
\begin{itemize}
    \item Subscription-based funding
    \item Callback pattern for async delivery
    \item Request confirmations for block reorg protection
\end{itemize}

This precompile enables synchronous, gas-efficient verification.

\subsection{RANDAO}

Ethereum's RANDAO accumulates validator contributions:
\[
\text{randao}_n = \text{xor}(\text{randao}_{n-1}, \text{reveal}_n)
\]

RANDAO is biasable by the last revealer; VRF eliminates this.

\subsection{drand}

drand provides distributed randomness beacons using BLS threshold signatures. The BLS-VRF suite enables similar functionality on-chain.

\section{Conclusion}

The VRF precompile enables efficient, verifiable randomness generation within the EVM. By providing native support for ECVRF verification, applications can obtain provably fair random numbers at a fraction of current costs. The Chainlink-compatible suite ensures interoperability with existing infrastructure, while the BLS-VRF extension enables threshold randomness for maximum decentralization. This capability unlocks secure randomness for gaming, NFT reveals, validator selection, and any application requiring trustless entropy.

\section*{Acknowledgments}

We thank the Chainlink team for their VRF specification, the IETF CFRG for RFC 9381 standardization, and the drand team for distributed randomness research.

\bibliographystyle{plain}
\begin{thebibliography}{9}

\bibitem{vrf}
Micali, S., Rabin, M., Vadhan, S.,
\textit{Verifiable Random Functions},
FOCS 1999.

\bibitem{rfc9381}
IETF,
\textit{RFC 9381: Verifiable Random Functions (VRFs)},
2023.

\bibitem{chainlink}
Chainlink,
\textit{Chainlink VRF v2 Documentation},
2022.

\bibitem{dvrf}
Galindo, D., et al.,
\textit{A Practical Distributed Verifiable Random Function},
2020.

\bibitem{ecvrf}
Papadopoulos, D., et al.,
\textit{Making NSEC5 Practical for DNSSEC},
IACR ePrint 2017.

\end{thebibliography}

\end{document}
